\par
\SourceLeaf{272}
Die Vergleichung der Zahlen im Thesaurus bei den 21 Bögen von $15^{0}38'20''$ bis $15^{0}41'40''$ mit den schärfer berechneten hat ergeben (IV):
\begin{center}\begin{tabular}{r r r r}
& Sin. & Cos. & Tang.\\
0&4&12&1\\
1&9&8&8\\
2&6&1&6\\
3&2&&4\\
4&&&2
\end{tabular}\end{center}
Die Fälle, wo die grössten Abweichungen vorkommen, sind:
\[
\begin{array}{l l r r r}
15^{0}40'20'',&\text{wo die Correctionen}&+3,&-1,&+4\\
15\quad41\quad30,&\text{wo folgende}&+3,&0,&+4
\end{array}
\]
an die Logarithmen des Sinus, des Cosinus und der Tangente angebracht werden müssen.

\Needspace{11\baselineskip}Ebenso hat die Vergleichung der \textsc{Vega}'schen Zahlen bei den 93 Bögen, von $1^{0}19'52''$ bis $1^{0}21'24''$ mit schärferer Rechnung folgende ergeben (V):
\begin{center}\begin{tabular}{r r r r}
& Sin. & Cos. & Tang.\\
0&38&30&17\\
1&39&56&41\\
2&15&7&22\\
3&1&&11\\
4&&&2
\end{tabular}\end{center}
Die grössten Abweichungen finden Statt bei den Bögen $1^{0}20'10''$ und $1^{0}20'15''$; und die nöthigen Correctionen betragen bei ersterm $-2,+1,-4$, bei dem andern $-2,+2,-4$.

Die Gruppen IV und V haben nicht gleichen Ursprung, da die Zahlen des Thesaurus zu der Gruppe IV aus \textsc{Vlacq}'s Trigonometria Artificialis genommen sind, die andern hingegen zu den neu hinzugekommenen gehören, welche unter \textsc{Vega}'s Leitung und Aufsicht von dem Lieutenant \textsc{Dorfmund} berechnet sind. Nach obigen Abrissen erscheinen die letztern wie etwas weniger ungenau als die erstern, wiewohl die Gruppe IV zu wenig zahlreich ist, um ein sicheres Urtheil zu begründen. Jedenfalls folgt aus dem Zusammenwerfen beider Gruppen eine Schätzung
\SourceLeaf{273}
für die Totalungenauigkeit der Tafeln, die eher etwas zu günstig sein wird, als umgekehrt. Aus dieser Vereinigung folgt, für 114 Bögen (VI):
\begin{center}\begin{tabular}{r r r r}
& Sin. & Cos. & Tang.\\
0&42&42&18\\
1&48&64&49\\
2&21&8&28\\
3&3&&15\\
4&&&4
\end{tabular}\end{center}
Es sind also hier unter 342 Logarithmen nur 102. die keiner Verbesserung bedürfen, gegen 240 ungenaue. Nach diesem Verhältniss würde man unter den 68038 irrationalen Logarithmen 47746 ungenaue erwarten können.

Die Summe der Quadrate der Abweichungen findet sich für die Sinus 159, für die Cosinus 96, für die Tangenten 360. Als mittlern Fehler mag man also annehmen für die Sinus $1{,}18$, für die Cosinus $0{,}92$, für die Tangenten $1{,}78$. Man kann hienach nicht zweifeln, dass die Tangenten-Logarithmen die abgeleiteten sind.

\begin{samepage}
Dass die Zahlen der Cosinuscolumne weniger ungenau sind, als die der Sinuscolumne, rührt wohl ohne Zweifel wenigstens theilweise, daher, dass bei den erstern die zur Ausfüllung erforderlichen Interpolationsmethoden einfacher ausfallen, möglicherweise können indess noch andere Ursachen mitgewirkt haben, worüber sich nur unsichere Vermuthungen aufstellen lassen würden.
\begin{center}\bigskip\rule{.3\linewidth}{.4pt}\end{center}
\end{samepage}
